
\def \Subject {برچسب‌گذاری دنباله}
\def \Session { یک}
% \setcounter{chapter}{\Session-1}
\renewcommand{\chaptername}{سوال}
\chapter{\Subject}
\section{بررسی پایگاه داده}
\begin{enumerate}
\item 

یکی از سطرهای مجموعه داده 
\LTRfootnote{Dataset}
در جدول 
\ref{tab:q1_datasample}
آمده است. می‌توان دید علاوه بر 
\lr{NER Tag}
مجموعه داده دو ستون دیگر به اسم‌های 
\lr{pos tag}
و 
\lr{chunk tag }
نیز دارد. 
\item 
\lr{pos tag}
عملا نقش گرامری هر توکن را مشخص می‌کند. برای مثال
\lr{NNP}
یک اسم مشخصه 
\LTRfootnote{Proper noun}
است. اهمیت این برچسب در 
\lr{NER}
از این رو است که ما به دنبال تشخیص کلاس اسامی (و به طور خاص اسامی مشخصه) هستیم. پس می‌توان افعال را درنظر نگرفت و اگر یک کلمه فعل بود از آن صرفنظر کرد. 

\lr{chunk tag}
عملا تعمیم 
\lr{pos tag}
است. یعنی سعی می‌کند جمله را بجای در سطح توکن در سطح عبارت
\LTRfootnote{phrase}
بسنجد. و توکن‌ها را باهم دسته می‌کند. 

  \begin{table}[ht]
	\centering
	\caption{سطر اول مجموعه داده آموزش}
	\label{tab:q1_datasample}
	\begin{latin}
\footnotesize

\begin{tabular}{llllllllll}
	\toprule
	\textbf{Token}  & EU & rejects & German & call & to & boycott & British & lamb & . \\
	\midrule
	\textbf{Pos Tag} & NNP & VBZ & JJ & NN & TO & VB & JJ & NN & . \\
	\textbf{Chunk Tag}& B-NP & B-VP & B-NP & I-NP & B-VP & I-VP & B-NP & I-NP & O \\
	\textbf{NER Tag} & B-ORG & O & B-MISC & O & O & O & B-MISC & O & O \\
	\bottomrule
\end{tabular}
	\end{latin}
	
	
\end{table}

\end{enumerate}
	
